% ============================================================================
%  capa.tex — capa, ficha tecnica e controle de versoes
% ============================================================================

\begin{titlepage}
\thispagestyle{empty}
\singlespacing
\setlength{\parskip}{0pt}
\begin{tikzpicture}[remember picture, overlay]
  % Faixa superior canônica — Trindade BLUNOORIA (Estrutura, Inteligência, Interface)
  \fill[unooNavy] (current page.north west) rectangle
        ([yshift=-12mm]current page.north east);
  \fill[unooLavender] ([yshift=-12mm]current page.north west) rectangle
        ([yshift=-13.8mm]current page.north east);
  \fill[unooBlue] ([yshift=-13.8mm]current page.north west) rectangle
        ([yshift=-15.0mm]current page.north east);

  % Marca d'água institucional — Símbolo Infinito Orgânico da Unooria no canto inferior
  \node[anchor=south east, inner sep=0, opacity=0.06] at ([xshift=-10mm, yshift=15mm]current page.south east) {
    \includegraphics[width=90mm]{figuras/marca/simbolo_azul_profundo.png}
  };

  % Barra lateral de estrutura e inteligência
  \fill[unooNavy] ([xshift=0mm, yshift=45mm]current page.south west) rectangle
        ([xshift=3.5mm, yshift=160mm]current page.south west);
  \fill[unooLavender] ([xshift=3.5mm, yshift=55mm]current page.south west) rectangle
        ([xshift=4.5mm, yshift=150mm]current page.south west);
\end{tikzpicture}

\vspace*{12mm}

% Logotipo oficial Unooria (Manual de Marca — Safespace e proporções originais)
\noindent\includegraphics[width=60mm]{figuras/marca/logo_azul_profundo.png}\par
\vspace{2mm}
{\footnotesize\color{unooLavenderSoft}\textbf{TECNOLOGIA EDUCACIONAL B2G}\quad\textperiodcentered\quad{\color{unooMuted}PLATAFORMA PÚBLICA INTEGRADA}\par}

\vspace{14mm}

{\footnotesize\bfseries\color{unooLavenderSoft}%
 DOCUMENTAÇÃO TÉCNICA DE ARQUITETURA DE SOFTWARE, ARQUITETURA DE NUVEM E CÓDIGO}

\vspace{3.5mm}

{\color{unooBlue}\rule{32mm}{2.6pt}}\hspace{2mm}{\color{unooLavender}\rule{12mm}{2.6pt}}

\vspace{8mm}

{\fontsize{32}{38}\selectfont\bfseries\color{unooNavy}%
 Plataforma Unooria\par}

\vspace{4mm}

{\fontsize{15}{21}\selectfont\bfseries\color{unooLavenderSoft}%
 Hub multi-tenant, 48 experiências imersivas e governança por escola\par}

\vspace{6mm}

{\fontsize{13.5}{19.5}\selectfont\color{unooInkSoft}%
 \textit{Monorepo} de uma plataforma de educação básica para redes públicas de
 ensino: um \emph{hub} em Next.js que embeda experiências estáticas versionadas
 em CDN, uma API NestJS com isolamento lógico por escola em \textit{Row-Level
 Security}, telemetria em dois canais e uma camada de IA generativa com quota
 por perfil --- uma \emph{stack} de nuvem por contrato, cem escolas por
 \emph{stack}.\par}

\vspace{11mm}

{\color{unooGrid}\rule{\linewidth}{0.8pt}}

\vspace{5mm}

\begin{tabularx}{\linewidth}{@{}L{30mm}@{\hspace{4mm}}X@{}}
  {\footnotesize\color{unooMuted}DOCUMENTO}      & {\footnotesize\bfseries\color{unooInk}\docid{} --- Documentação Técnica Consolidada}\\[3pt]
  {\footnotesize\color{unooMuted}VERSÃO}         & {\footnotesize\bfseries\color{unooInk}\docversao{} (Modelagem UML 2.5.1)}\\[3pt]
  {\footnotesize\color{unooMuted}DATA}           & {\footnotesize\bfseries\color{unooInk}09 de setembro de 2026}\\[3pt]
  {\footnotesize\color{unooMuted}AUTOR}          & {\footnotesize\bfseries\color{unooInk}Lucas Salomão}\\[3pt]
  {\footnotesize\color{unooMuted}REPOSITÓRIO}    & {\footnotesize\bfseries\color{unooInk}\texttt{plataforma-unooria} (\textit{monorepo}, tronco \texttt{main})}\\[3pt]
  {\footnotesize\color{unooMuted}CLASSIFICAÇÃO}  & {\footnotesize\bfseries\color{unooInk}Documento técnico de projeto --- uso interno}\\
\end{tabularx}

\vspace{5mm}

{\color{unooGrid}\rule{\linewidth}{0.8pt}}

\vspace{6mm}

{\scriptsize\color{unooMuted}%
 Escopo do documento: contexto de negócio, requisitos, casos de uso,
 arquitetura de \textit{software} nos quatro níveis do modelo C4, arquitetura de
 nuvem e implantação em GCP, modelo de \textit{multi-tenancy} e segurança,
 modelo de dados, contratos de integração, telemetria, camada de IA, catálogo de
 experiências, especificação da API, \textit{design system}, engenharia de
 entrega, qualidade, registro de decisões arquiteturais, riscos e evolução.
 Diagramas vetoriais em TikZ; fontes Mermaid equivalentes no Apêndice~A.\par}

\end{titlepage}

% ---------------------------------------------------------------------------
%  Ficha tecnica e controle de versoes
% ---------------------------------------------------------------------------
\chapter*{Ficha técnica do documento}
\addcontentsline{toc}{chapter}{Ficha técnica do documento}
\markboth{Ficha técnica do documento}{}

\section*{Identificação}

\begin{tabularx}{\linewidth}{@{}L{36mm}X@{}}
  \toprule
  \textbf{Campo} & \textbf{Conteúdo}\\
  \midrule
  Título & Plataforma Unooria --- Documentação Técnica Consolidada de
           Arquitetura de \textit{Software}, Arquitetura de Nuvem e Código\\
  \zebra Identificador & \docid\\
  Versão & \docversao\\
  \zebra Data de emissão & 27 de agosto de 2026\\
  Autor & Lucas Salomão (\url{lucas.salomao@unooria.com})\\
  \zebra Repositório & \arq{plataforma-unooria} --- \textit{monorepo} pnpm + Turborepo\\
  Tronco & \arq{main} (histórico linear; uma \textit{branch} por experiência)\\
  \zebra Linguagem principal & TypeScript 5.7 (\textit{strict}), Node.js 20\\
  Prazo de produção & novembro de 2026\\
  \zebra Estado & Em desenvolvimento --- M0--M6 concluídos; \textit{stack} de
           \textit{dev} em GCP operacional\\
  Repositório de infra & Separado (Terraform). Contrato em \arq{docs/INFRA.md}\\
  \bottomrule
\end{tabularx}

\section*{Escopo e propósito}

Este documento consolida, em fonte única, o que estava distribuído entre
\arq{AGENTS.md} (as convenções e regras invioláveis do repositório),
\arq{README.md}, \arq{docs/ARQUITETURA.md}, \arq{docs/DECISOES.md},
\arq{docs/INFRA.md}, \arq{docs/INFRA-IDENTITY.md} e os comentários de decisão
que vivem no próprio código-fonte --- que, neste projeto, é onde as decisões
mais finas estão registradas.

Acrescenta o material que não existia em forma escrita: especificação de
requisitos funcionais e não funcionais rastreáveis, catálogo de casos de uso,
modelo C4 completo nos quatro níveis para o \textit{software} \emph{e} para a
nuvem, diagramas de sequência dos contratos de integração, dicionário de dados,
especificação da API rota por rota e o registro formal das decisões
arquiteturais (ADR) com suas consequências.

\begin{nota}
Este documento é \textbf{normativo} para o projeto: quando ele divergir de uma
suposição, ele vence. Quando divergir do código, o código vence e o documento
está desatualizado --- e atualizá-lo faz parte da \emph{Definition of Done} de
qualquer tarefa (Capítulo~\ref{cap:qualidade}).
\end{nota}

\section*{Histórico de revisões}

\begin{tabularx}{\linewidth}{@{}L{14mm}L{22mm}L{30mm}X@{}}
  \toprule
  \headrow \thd{Versão} & \thd{Data} & \thd{Autor} & \thd{Alterações}\\
  \midrule
  0.1 & 22/08/2026 & L. Salomão & Estrutura do \textit{monorepo}; arquitetura e
       ADRs iniciais em \arq{docs/}.\\
  \zebra 0.2 & 25/08/2026 & L. Salomão & Migração para o \textit{monorepo}
       mergeada na \arq{main} (26 \textit{commits}, \textit{fast-forward});
       \textit{tenancy} com RLS, exp08ef1b e exp04ef1b.\\
  0.3 & 26/08/2026 & L. Salomão & app-70 (Vila do Lume) e app-91 (bem-estar);
       decisão de retenção sem prazo declarado.\\
  \zebra \textbf{1.0} & \textbf{27/08/2026} & \textbf{L. Salomão} &
       \textbf{Documentação técnica consolidada: requisitos, casos de uso,
       modelo C4 de \textit{software} e de nuvem, \textit{multi-tenancy},
       modelo de dados, contratos, telemetria particionada, IA, API,
       \textit{design system}, entrega, qualidade, ADRs e bibliografia em
       ABNT.}\\
  \zebra 1.1 & 01/09/2026 & L. Salomão & Capítulo de observabilidade: estado
       atual por camada, sinais, comparação de ferramentas, SLOs e alertas,
       privacidade nos sinais, custo por contrato e roteiro em fases;
       \adrr{ADR-57} e \adrr{ADR-58} (propostas); \cd{R-05} reclassificado.\\
  \textbf{1.2} & 09/09/2026 & L. Salomão & Revisão técnica integral e consolidação
       multi-agente: atualização do modelo de 6 papéis (\adrr{ADR-80}), catálogo de 49 tabelas
       e 30 migrações Prisma, 85 ADRs consolidadas, novo Capítulo~\ref{cap:pendencias} de
       pendências rastreadas, adequação integral da identidade visual e layout ao
       Manual de Marca oficial (\textit{Brandbook Unooria}).\\
  \textbf{1.3} & 09/09/2026 & L. Salomão & Terceiro eixo de acesso
       (\emph{disponibilidade}) e painel de recursos da Unooria: \adrr{ADR-87} e
       \adrr{ADR-88}, tabelas \cd{disponibilidades} e
       \cd{experiencias\_versoes}, campo \cd{produto} no manifesto de
       experiência.\\
  \zebra \textbf{1.4} & 27/09/2026 & L. Salomão & Modelagem UML: os catorze tipos
       de diagrama da UML 2.5.1 no Capítulo~\ref{cap:arqsw}
       (Seção~\ref{sec:uml}), com o perfil da plataforma, conferidos contra o
       código, e o par Mermaid de cada um (\adrr{ADR-122}); casos de uso com os
       seis papéis; correções que a conferência expôs (grafo do
       \emph{monorepo}, quota na camada de IA). Acumula as \adrr{ADR-89} a
       \adrr{ADR-121}, registradas desde a 1.3 sem revisão de versão.\\
  \bottomrule
\end{tabularx}

\section*{Documentos de referência do projeto}

\begin{tabularx}{\linewidth}{@{}L{45mm}X@{}}
  \toprule
  \headrow \thd{Artefato} & \thd{Conteúdo}\\
  \midrule
  \arq{AGENTS.md} & Convenções, regras invioláveis de \textit{tenancy} e
       segurança, contratos de integração e \emph{Definition of Done}. É o
       documento de governança do repositório.\\
  \zebra \arq{README.md} & Instalação, os três modos de execução local e
       \textit{smoke} visual.\\
  \arq{Manual de Marca} & Manual oficial de identidade visual (\textit{Brandbook Unooria}):
       paleta canônica BLUNOORIA (Azul Profundo, Lavanda, Azul Claro), proporções e safespace
       do logotipo e do símbolo infinito orgânico.\\
  \zebra \arq{docs/ARQUITETURA.md} & Visão arquitetural resumida por contrato.\\
  \arq{docs/DECISOES.md} & Registro resumido de decisões (base do
       Capítulo~\ref{cap:adr}).\\
  \zebra \arq{docs/INFRA.md} & Contrato entre este repositório e o de Terraform:
       nomes de recursos, variáveis e o que o CI publica.\\
  \arq{docs/INFRA-IDENTITY.md} & Contrato do gerenciador central de
       identidade (WSO2 IS), que é único para vários contratos.\\
  \zebra \arq{apps/api/prisma/schema.prisma} & Modelo de dados anotado, com o porquê
       de cada decisão de esquema.\\
  \arq{packages/sdk/src/} & Contrato \textit{postMessage} e cliente de
       telemetria --- a fronteira entre hub e experiência.\\
  \zebra \arq{.env.example} & Catálogo completo de variáveis de ambiente
       (Apêndice~\ref{ap:envs}).\\
  \bottomrule
\end{tabularx}
